\documentclass[a4paper,11pt]{article}
\usepackage[utf8]{inputenc}
\usepackage[T1]{fontenc}
\usepackage[french]{babel}
\usepackage[left=2cm,right=2cm,top=2cm,bottom=2cm]{geometry}
\usepackage{xcolor}
\usepackage{hyperref}
\usepackage{fontawesome5}
\usepackage{parskip}

% --- Couleurs Accenture ---
\definecolor{primary}{RGB}{117, 59, 189}
\hypersetup{colorlinks=true, urlcolor=primary}

\begin{document}

% --- EXPÉDITEUR ---
\noindent
\begin{minipage}[t]{0.5\textwidth}
    \textbf{Loïc Aron MBASSI EWOLO}\\
    Élève-Ingénieur Bac+5 -- Double diplôme ENSEA\\[3pt]
    \faMapMarker\ Cergy, France\\
    \faPhone\ (+33) 7 59 52 33 98\\
    \faEnvelope\ \href{mailto:loicaron.mbassiewolo@ensea.fr}{loicaron.mbassiewolo@ensea.fr}
\end{minipage}
\hfill
% --- DATE ---
\begin{minipage}[t]{0.4\textwidth}
    \raggedleft
    Cergy, le 5 février 2026
\end{minipage}

\vspace{1.2cm}

% --- DESTINATAIRE ---
\hfill
\begin{minipage}[t]{0.5\textwidth}
    \raggedleft
    \textbf{Accenture France}\\
    Service Recrutement\\
    Clermont-Ferrand (63), France
\end{minipage}

\vspace{1cm}

\noindent\textbf{Objet :} Candidature au stage Développeur IA (Clermont-Ferrand)

\vspace{0.8cm}

Madame, Monsieur,

Actuellement élève-ingénieur en double diplôme à l'ENSEA et à l'École Polytechnique de Yaoundé, spécialisé en Intelligence Artificielle et Machine Learning, je suis à la recherche d'un stage de fin d'études à partir d'avril 2026. Votre offre de stage sur l'IA générative appliquée à l'amélioration de l'accès à l'information correspond parfaitement à mon projet professionnel et à mes compétences.

Mon expérience la plus pertinente pour ce stage est la conception d'un système multi-agents agricole utilisant RAG (Retrieval-Augmented Generation) et LangChain. J'y ai orchestré quatre agents IA collaboratifs via Google ADK, implémenté l'enrichissement contextuel depuis une base de connaissances, et conçu des prompts adaptés aux différents cas d'usage. Cette expérience m'a confronté aux enjeux que vous mentionnez : structuration des données, pertinence des réponses générées et évaluation des performances.

Par ailleurs, mon projet Women Safe (détection de contenus sexistes avec BERT et GPT-2) m'a sensibilisé aux questions de Responsible AI et de biais éthiques, des principes essentiels chez Accenture. Mon stage de recherche au MILA (Institut Québécois d'IA) m'a permis de développer une rigueur scientifique dans l'évaluation des modèles et la reproduction d'expériences.

Je maîtrise Python, PyTorch, LangChain et les techniques de prompting. Je suis capable de contribuer à l'analyse des besoins fonctionnels, à l'élaboration de modèles de requêtes pertinents et au développement d'un prototype respectant les bonnes pratiques de qualité et de sécurité.

Je serais honoré de rejoindre Accenture pour contribuer à ce projet innovant et apprendre auprès de vos experts. Je reste à votre disposition pour un entretien.

Je vous prie d'agréer, Madame, Monsieur, l'expression de mes salutations distinguées.

\vspace{0.8cm}

\textbf{Loïc Aron MBASSI EWOLO}\\
\textit{Élève-Ingénieur ENSEA / Polytechnique Yaoundé}

\end{document}
